\apendice{Termo de Assentimento Livre e Esclarecido (TALE)}
\label{ap:tale}

\begin{center}
    \textbf{UNIVERSIDADE DO ESTADO DA BAHIA (UNEB)}\\[0.2cm]
    \textbf{DEPARTAMENTO DE CIÊNCIAS EXATAS E DA TERRA (DCET)}\\[0.2cm]
    \textbf{CURSO DE GRADUAÇÃO EM SISTEMAS DE INFORMAÇÃO}
\end{center}

\noindent Esta pesquisa seguirá os critérios da ética em pesquisa com seres humanos, conforme a Resolução CNS nº 466/12 e a Resolução CNS nº 510/16.

\section{Convite}
Olá! Você está sendo convidado(a) a participar da pesquisa \imprimirtitulo, que está sendo realizada pelo pesquisador \imprimirautor\ e pela professora \imprimirorientador, da Universidade do Estado da Bahia (UNEB). Seus pais ou responsáveis já permitiram que você participe e assinaram um documento autorizando, mas a decisão final de participar ou não é sua.

\section{Por que estamos fazendo esta pesquisa?}
Queremos entender melhor se jogar um jogo de perguntas de Matemática do \gls{enem} ajuda você a aprender de um jeito divertido, a resolver problemas com mais calma e a gostar mais dessa matéria. A nossa ideia é descobrir se o jogo ajuda estudantes como você a se preparar para a prova.

\section{O que você vai fazer?}
Se você aceitar participar, nós vamos usar com você uma plataforma de perguntas de Matemática do \gls{enem}, no computador, com pontos, medalhas e desafios. Depois, vamos fazer algumas perguntinhas rápidas sobre o que você achou da experiência, que você responderá em bem pouquinho tempo. Nós vamos ficar junto de você durante a atividade e explicar tudo o que for preciso.

\section{E se você ficar cansado(a) ou não quiser participar?}
Você só participa se quiser! Se você não quiser participar, não tem problema nenhum, e ninguém vai ficar chateado com você. Se você aceitar agora e depois quiser parar ou não responder a alguma pergunta, você pode avisar a qualquer momento, e nós paramos na mesma hora, sem nenhum prejuízo para você.

\section{Existe alguma coisa ruim que pode acontecer?}
Talvez você fique chateado(a) ou frustrado(a) se errar uma questão, ou um pouco cansado(a) se ficar muito tempo jogando. Para ajudar, no jogo o erro faz parte do aprendizado e você pode pedir dicas quando precisar; as atividades são curtinhas e podem ser interrompidas quando você quiser.

\section{O que pode ser bom para você?}
Você vai poder usar um jogo para estudar Matemática, conhecer o tipo de questão que aparece no \gls{enem} e praticar bastante. Você também vai nos ajudar a melhorar a plataforma para outros estudantes.

\section{Suas informações ficarão em segredo}
Ninguém da sua escola nem os seus amigos saberão o que você respondeu. Não vamos colocar o seu nome nos trabalhos: usaremos um nome inventado para proteger você. Quando os resultados da pesquisa forem publicados em jornais, revistas ou eventos científicos, ninguém saberá quem participou, e você também poderá ter acesso a esses resultados.

\section{E se você tiver dúvidas?}
Se você ou seus pais (ou responsáveis) tiverem qualquer dúvida, vocês podem procurar:

\begin{itemize}
    \item Pesquisador discente (orientando): \imprimirautor\ -- Telefone: \rule{3cm}{0.4pt}\ E-mail: \rule{4cm}{0.4pt}
    \item Professora responsável (orientadora): \imprimirorientador\ -- Telefone: \rule{3cm}{0.4pt}\ E-mail: \rule{4cm}{0.4pt}
    \item Comitê de Ética em Pesquisa com Seres Humanos da UNEB (CEP/UNEB): Edifício Civil Empresarial, Rua Dr. José Peroba, 251, 7º andar -- Stiep, Salvador -- BA, CEP 41770-235. Telefone: (71) 3198-5856. E-mail: cepuneb@uneb.br
\end{itemize}

\section{Seus pais ou responsáveis já autorizaram}
Nós já conversamos com os seus pais (ou responsáveis) e eles autorizaram a sua participação por escrito, assinando o Termo de Consentimento Livre e Esclarecido (apêndice \ref{ap:tcle-responsavel}), mas a decisão final de participar ou não é sua.

\section*{Declaração de assentimento}
Eu, \rule{6cm}{0.4pt}, aceito participar da pesquisa \imprimirtitulo. Entendi os objetivos, as coisas boas e as coisas ruins que podem acontecer. Entendi que posso dizer ``sim'' e participar, mas que, a qualquer momento, posso dizer ``não'' e desistir, sem nenhum problema. Os pesquisadores tiraram minhas dúvidas e conversaram com os meus responsáveis. Recebi uma cópia deste termo de assentimento, li e concordo em participar da pesquisa.

\vspace{0.5cm}

\noindent Salvador -- BA, \rule{2cm}{0.4pt} de \rule{3cm}{0.4pt} de 20\rule{1cm}{0.4pt}.

\vspace{1.5cm}

\noindent\rule{9cm}{0.4pt}\\
Assinatura (ou nome escrito) do(a) participante da pesquisa

\vspace{1.5cm}

\noindent\rule{9cm}{0.4pt}\\
Assinatura do(a) pesquisador(a) responsável